
\documentclass{article}
\usepackage{colortbl}
\usepackage{makecell}
\usepackage{multirow}
\usepackage{supertabular}

\begin{document}

\newcounter{utterance}

\centering \large Interaction Transcript for game `function\_detective', experiment `two\_numbers\_passive\_labeled\_pairs', episode 243 with mistral{-}large{-}3{-}azure.
\vspace{24pt}

{ \footnotesize  \setcounter{utterance}{1}
\setlength{\tabcolsep}{0pt}
\begin{supertabular}{c@{$\;$}|p{.15\linewidth}@{}p{.15\linewidth}p{.15\linewidth}p{.15\linewidth}p{.15\linewidth}p{.15\linewidth}}
    \# & \multicolumn{2}{c}{Player} && \multicolumn{2}{c}{Game Master} \\
    \hline

    \theutterance \stepcounter{utterance}  
    & & & \multicolumn{4}{p{0.6\linewidth}}{
        \cellcolor[rgb]{0.9,0.9,0.9}{
            \makecell[{{p{\linewidth}}}]{
                \texttt{\tiny{[P1$\langle$GM]}}
                \texttt{I have a hidden Python function for you to figure out.} \\
\\ 
\texttt{Your job is to infer what it does and then submit a correct Python implementation.} \\
\\ 
\texttt{Function signature:} \\
\texttt{(a: int, b: int) {-}> int} \\
\\ 
\texttt{You have 15 turns total.} \\
\\ 
\texttt{I reveal candidate input{-}output pairs one at a time and tell you whether each one is correct.} \\
\\ 
\texttt{You must make your best SOLVE attempt no later than the second{-}to{-}last turn.} \\
\texttt{Do not wait until the very end.} \\
\\ 
\texttt{You can use these actions:} \\
\texttt{{-} JSON format described at the end of the prompt} \\
\texttt{{-} SOLVE: ```python ... ```} \\
\\ 
\texttt{Do not add import statements in your guess.} \\
\texttt{Anything you might need is already available, so just write the function body in hypothesis and return the answer.} \\
\\ 
\texttt{Preloaded examples:} \\
\texttt{1. (((99, 1000), 9), True)} \\
\\ 
\texttt{Scoring:} \\
\texttt{{-} correctness is required for any non{-}zero score} \\
\texttt{{-} in interactive modes, solving with fewer information requests is better} \\
\\ 
\texttt{Input formatting reminders:} \\
\texttt{{-} strings must be quoted} \\
\texttt{{-} booleans must be True/False} \\
\texttt{{-} use commas between multiple values} \\
\texttt{{-} do not write any extra action tags in the same message} \\
\\ 
\texttt{Format example:} \\
\texttt{You write:} \\
\texttt{\{'current\_mode': <probing/confirming/don't know; choose whichever best describes your current state in the task>, 'input\_rationale': <brief description of why you chose this hypothesis>, 'hypothesis': <either the exact string "unknown" if you have no guess yet OR your current best guess as runnable Python matching the given signature, with a one{-}line docstring stating what it does. Never write placeholder code, if you cannot commit to a guess, use "unknown". Match the given signature.} \\
\texttt{e.g.} \\
\texttt{def f(x):} \\
\texttt{	"""Returns the absolute value of x."""} \\
\texttt{	return x if x >= 0 else {-}x} \\
\\ 
\texttt{Use "unknown" if you have no guess.>, 'confidence\_score': <score between 1{-}100. 100 means highest confidence in your hypothesis>\}} \\
\\ 
\texttt{Output only one valid json per turn.} \\
            }
        }
    }
     \\ \\

    \theutterance \stepcounter{utterance}  
    & \multicolumn{4}{p{0.6\linewidth}}{
        \cellcolor[rgb]{0.561,0.737,0.561}{
            \makecell[{{p{\linewidth}}}]{
                \texttt{\tiny{[P1$\rangle$GM]}}
                \texttt{```json} \\
\texttt{\{} \\
\texttt{    "current\_mode": "probing",} \\
\texttt{    "input\_rationale": "The preloaded example ((99, 1000), 9) suggests that the function might be returning the difference between the two numbers, the minimum, or some other simple arithmetic operation. Starting with a simple test case to check if it's the difference (a {-} b or b {-} a).",} \\
\texttt{    "hypothesis": "unknown",} \\
\texttt{    "confidence\_score": 1} \\
\texttt{\}} \\
\texttt{```} \\
            }
        }
    }
    & & \\ \\

    \theutterance \stepcounter{utterance}  
    & & & \multicolumn{4}{p{0.6\linewidth}}{
        \cellcolor[rgb]{0.9,0.9,0.9}{
            \makecell[{{p{\linewidth}}}]{
                \texttt{\tiny{[P1$\langle$GM]}}
                \texttt{OUTPUT: ((({-}99, 0), 19), False)} \\
            }
        }
    }
     \\ \\

    \theutterance \stepcounter{utterance}  
    & \multicolumn{4}{p{0.6\linewidth}}{
        \cellcolor[rgb]{0.561,0.737,0.561}{
            \makecell[{{p{\linewidth}}}]{
                \texttt{\tiny{[P1$\rangle$GM]}}
                \texttt{```json} \\
\texttt{\{} \\
\texttt{    "current\_mode": "probing",} \\
\texttt{    "input\_rationale": "The first example ((99, 1000), 9) suggested a possible difference (1000 {-} 99 = 901, which doesn't match), or perhaps the number of digits in the difference. The new example (({-}99, 0), 19) is False, which rules out simple differences. The output 19 for ({-}99, 0) suggests it might be related to the absolute value of the numbers or their digit counts. Hypothesizing that it could be the sum of the absolute values of the digits of a and b.",} \\
\texttt{    "hypothesis": "def f(a: int, b: int) {-}> int:\textbackslash{}n    \textbackslash{}"\textbackslash{}"\textbackslash{}"Returns the sum of the absolute values of the digits of a and b.\textbackslash{}"\textbackslash{}"\textbackslash{}"\textbackslash{}n    def digit\_sum(x):\textbackslash{}n        return sum(int(d) for d in str(abs(x)))\textbackslash{}n    return digit\_sum(a) + digit\_sum(b)",} \\
\texttt{    "confidence\_score": 60} \\
\texttt{\}} \\
\texttt{```} \\
            }
        }
    }
    & & \\ \\

    \theutterance \stepcounter{utterance}  
    & & & \multicolumn{4}{p{0.6\linewidth}}{
        \cellcolor[rgb]{0.9,0.9,0.9}{
            \makecell[{{p{\linewidth}}}]{
                \texttt{\tiny{[P1$\langle$GM]}}
                \texttt{OUTPUT: (((2, {-}2), 4), True)} \\
            }
        }
    }
     \\ \\

    \theutterance \stepcounter{utterance}  
    & \multicolumn{4}{p{0.6\linewidth}}{
        \cellcolor[rgb]{0.561,0.737,0.561}{
            \makecell[{{p{\linewidth}}}]{
                \texttt{\tiny{[P1$\rangle$GM]}}
                \texttt{```json} \\
\texttt{\{} \\
\texttt{    "current\_mode": "confirming",} \\
\texttt{    "input\_rationale": "The previous hypothesis (sum of absolute digit values) works for the new example ((2, {-}2), 4): 2 + 2 = 4. The first example ((99, 1000), 9) also fits: 9+9 + 1+0+0+0 = 18, but the output was 9, so this hypothesis is incorrect. Re{-}evaluating. The output 9 for (99, 1000) could be the number of digits in the smaller number (99 has 2 digits, 1000 has 4, but 2+4=6 ≠ 9). Alternatively, it might be the sum of the digits of the smaller number (9+9=18 ≠ 9). Another possibility: the function returns the number of digits in the product of a and b. For (99, 1000): 99*1000=99000 (5 digits), but output is 9. Not matching. New hypothesis: the function returns the sum of the digits of a and b, but only for the smaller number in absolute value. For (99, 1000): min(99, 1000) is 99, sum of} \\
            }
        }
    }
    & & \\ \\

    \theutterance \stepcounter{utterance}  
    & & & \multicolumn{4}{p{0.6\linewidth}}{
        \cellcolor[rgb]{0.9,0.9,0.9}{
            \makecell[{{p{\linewidth}}}]{
                \texttt{\tiny{[P1$\langle$GM]}}
                \texttt{Game Violation: ParseError: Response is not valid JSON.} \\
\texttt{Output ONLY:} \\
\texttt{One JSON with confidence\_score, current\_mode, input\_rationale, and hypothesis keys ...} \\
\texttt{OR} \\
\texttt{SOLVE: ```python ... ```} \\
            }
        }
    }
     \\ \\

    \theutterance \stepcounter{utterance}  
    & \multicolumn{4}{p{0.6\linewidth}}{
        \cellcolor[rgb]{0.561,0.737,0.561}{
            \makecell[{{p{\linewidth}}}]{
                \texttt{\tiny{[P1$\rangle$GM]}}
                \texttt{```json} \\
\texttt{\{} \\
\texttt{    "current\_mode": "confirming",} \\
\texttt{    "input\_rationale": "The sum of absolute digits hypothesis failed for the first example ((99, 1000), 9), but worked for ((2, {-}2), 4). The correct pattern might instead be the count of digits in the smaller absolute value. For (99, 1000): min(99, 1000) is 99 (2 digits), but output is 9, which doesn't match. Alternatively, it could be the sum of the digits of the smaller absolute value: 9+9=18 ≠ 9. New hypothesis: the function returns the number of digits in `a` multiplied by the number of digits in `b`. For (99, 1000): 2 * 4 = 8 ≠ 9. Not matching. Revisiting the first example: 99 and 1000. 99 + 1000 = 1099 (sum of digits: 1+0+9+9=19 ≠ 9). Another idea: the function returns the length of the string representation of the smaller absolute value. For (99, 1000): len('99') = 2 ≠ 9. Not matching. Final hypothesis: the function returns the sum of} \\
            }
        }
    }
    & & \\ \\

    \theutterance \stepcounter{utterance}  
    & & & \multicolumn{4}{p{0.6\linewidth}}{
        \cellcolor[rgb]{0.9,0.9,0.9}{
            \makecell[{{p{\linewidth}}}]{
                \texttt{\tiny{[P1$\langle$GM]}}
                \texttt{Game Violation: ParseError: Response is not valid JSON.} \\
\texttt{Output ONLY:} \\
\texttt{One JSON with confidence\_score, current\_mode, input\_rationale, and hypothesis keys ...} \\
\texttt{OR} \\
\texttt{SOLVE: ```python ... ```} \\
            }
        }
    }
     \\ \\

    \theutterance \stepcounter{utterance}  
    & \multicolumn{4}{p{0.6\linewidth}}{
        \cellcolor[rgb]{0.561,0.737,0.561}{
            \makecell[{{p{\linewidth}}}]{
                \texttt{\tiny{[P1$\rangle$GM]}}
                \texttt{```json} \\
\texttt{\{} \\
\texttt{    "current\_mode": "confirming",} \\
\texttt{    "input\_rationale": "The sum of absolute digits hypothesis worked for ((2, {-}2), 4) but failed for ((99, 1000), 9). The correct pattern might be the number of digits in the smaller number plus the number of digits in the larger number. For (99, 1000): 2 + 4 = 6 ≠ 9. Alternatively, it could be the sum of the digits of the smaller number multiplied by some factor. For (99, 1000): 9+9=18, 18/2=9, which matches the output. For (2, {-}2): 2+2=4, 4/1=4, which matches. Testing this hypothesis.",} \\
\texttt{    "hypothesis": "def f(a: int, b: int) {-}> int:\textbackslash{}n    \textbackslash{}"\textbackslash{}"\textbackslash{}"Returns the sum of the digits of the smaller absolute value.\textbackslash{}"\textbackslash{}"\textbackslash{}"\textbackslash{}n    def digit\_sum(x):\textbackslash{}n        return sum(int(d) for d in str(abs(x)))\textbackslash{}n    return digit\_sum(min(a, b, key=abs))",} \\
\texttt{    "confidence\_score": 80} \\
\texttt{\}} \\
\texttt{```} \\
            }
        }
    }
    & & \\ \\

    \theutterance \stepcounter{utterance}  
    & & & \multicolumn{4}{p{0.6\linewidth}}{
        \cellcolor[rgb]{0.9,0.9,0.9}{
            \makecell[{{p{\linewidth}}}]{
                \texttt{\tiny{[P1$\langle$GM]}}
                \texttt{OUTPUT: ((({-}1000, {-}1000), {-}1), False)} \\
            }
        }
    }
     \\ \\

    \theutterance \stepcounter{utterance}  
    & \multicolumn{4}{p{0.6\linewidth}}{
        \cellcolor[rgb]{0.561,0.737,0.561}{
            \makecell[{{p{\linewidth}}}]{
                \texttt{\tiny{[P1$\rangle$GM]}}
                \texttt{```json} \\
\texttt{\{} \\
\texttt{    "current\_mode": "probing",} \\
\texttt{    "input\_rationale": "The last hypothesis (sum of digits of the smaller absolute value) fails for (({-}1000, {-}1000), {-}1). This suggests the function does not return a negative value, and the output is not purely based on digit sums. The output {-}1 is likely a sentinel value or indicates something else. Given the examples, the function might return the number of digits in the GCD of a and b. For (99, 1000): GCD is 1 (1 digit) ≠ 9. For (2, {-}2): GCD is 2 (1 digit) ≠ 4. Not matching. Alternatively, it could be the count of set bits in the binary representation of the sum of a and b. For (99, 1000): 99+1000=1099 (binary 10001000011, 5 set bits) ≠ 9. For (2, {-}2): 2{-}2=0 (0 set bits) ≠ 4. Not matching. New hypothesis: the function returns the number of trailing zeros in the product of a and b. For (99, 1000): 99*1000=99000} \\
            }
        }
    }
    & & \\ \\

    \theutterance \stepcounter{utterance}  
    & & & \multicolumn{4}{p{0.6\linewidth}}{
        \cellcolor[rgb]{0.9,0.9,0.9}{
            \makecell[{{p{\linewidth}}}]{
                \texttt{\tiny{[P1$\langle$GM]}}
                \texttt{Game Violation: ParseError: Response is not valid JSON.} \\
\texttt{Output ONLY:} \\
\texttt{One JSON with confidence\_score, current\_mode, input\_rationale, and hypothesis keys ...} \\
\texttt{OR} \\
\texttt{SOLVE: ```python ... ```} \\
            }
        }
    }
     \\ \\

    \theutterance \stepcounter{utterance}  
    & \multicolumn{4}{p{0.6\linewidth}}{
        \cellcolor[rgb]{0.561,0.737,0.561}{
            \makecell[{{p{\linewidth}}}]{
                \texttt{\tiny{[P1$\rangle$GM]}}
                \texttt{```json} \\
\texttt{\{} \\
\texttt{    "current\_mode": "probing",} \\
\texttt{    "input\_rationale": "The previous hypotheses are failing. Let's re{-}examine the examples: ((99, 1000), 9), ((2, {-}2), 4), (({-}1000, {-}1000), {-}1). The output {-}1 for ({-}1000, {-}1000) suggests the function might return {-}1 when inputs are equal. For (99, 1000), 9 could be the number of letters in the English word for 99 ('ninety{-}nine' has 9 letters excluding hyphen). For (2, {-}2), 'two' has 3 letters, but output is 4, so this doesn't fit. Alternatively, 9 could be the count of letters in the Roman numeral for 99 (XCIX: 4 letters), but 4 ≠ 9. Not matching. New hypothesis: the function returns the count of prime factors (with multiplicity) of the smaller absolute value. For 99: 3*3*11 (3 factors) ≠ 9. For 2: 2 (1 factor) ≠ 4. Not matching. Final hypothesis: the function returns the sum of the digits of the smaller number multiplied by the number of digits in the larger number. For (99, 1} \\
            }
        }
    }
    & & \\ \\

    \theutterance \stepcounter{utterance}  
    & & & \multicolumn{4}{p{0.6\linewidth}}{
        \cellcolor[rgb]{0.9,0.9,0.9}{
            \makecell[{{p{\linewidth}}}]{
                \texttt{\tiny{[P1$\langle$GM]}}
                \texttt{Game Violation: ParseError: Response is not valid JSON.} \\
\texttt{Output ONLY:} \\
\texttt{One JSON with confidence\_score, current\_mode, input\_rationale, and hypothesis keys ...} \\
\texttt{OR} \\
\texttt{SOLVE: ```python ... ```} \\
            }
        }
    }
     \\ \\

    \theutterance \stepcounter{utterance}  
    & \multicolumn{4}{p{0.6\linewidth}}{
        \cellcolor[rgb]{0.561,0.737,0.561}{
            \makecell[{{p{\linewidth}}}]{
                \texttt{\tiny{[P1$\rangle$GM]}}
                \texttt{```json} \\
\texttt{\{} \\
\texttt{    "current\_mode": "probing",} \\
\texttt{    "input\_rationale": "After reviewing all examples, the pattern seems to be: if a == b, return {-}1. Otherwise, return the sum of the digits of the smaller absolute value multiplied by the number of digits in the larger absolute value. For (99, 1000): smaller=99 (digit sum=18), larger=1000 (4 digits), 18 * 0.5=9 (matches). For (2, {-}2): smaller=2 (digit sum=2), larger=2 (1 digit), 2 * 2=4 (matches). For ({-}1000, {-}1000): a == b, return {-}1 (matches). This seems consistent.",} \\
\texttt{    "hypothesis": "def f(a: int, b: int) {-}> int:\textbackslash{}n    \textbackslash{}"\textbackslash{}"\textbackslash{}"Returns {-}1 if a == b, else returns (sum of digits of min(abs(a), abs(b))) * (digits in max(abs(a), abs(b))) // 2.\textbackslash{}"\textbackslash{}"\textbackslash{}"\textbackslash{}n    if a == b:\textbackslash{}n        return {-}1\textbackslash{}n    def digit\_sum(x):\textbackslash{}n        return sum(int(d) for d in str(abs(x)))\textbackslash{}n    def digit\_count(x):\textbackslash{}n        return len(str(abs(x)))\textbackslash{}n    smaller = min(abs(a), abs(b))\textbackslash{}} \\
            }
        }
    }
    & & \\ \\

    \theutterance \stepcounter{utterance}  
    & & & \multicolumn{4}{p{0.6\linewidth}}{
        \cellcolor[rgb]{0.9,0.9,0.9}{
            \makecell[{{p{\linewidth}}}]{
                \texttt{\tiny{[P1$\langle$GM]}}
                \texttt{Game Violation: ParseError: Response is not valid JSON.} \\
\texttt{Output ONLY:} \\
\texttt{One JSON with confidence\_score, current\_mode, input\_rationale, and hypothesis keys ...} \\
\texttt{OR} \\
\texttt{SOLVE: ```python ... ```} \\
            }
        }
    }
     \\ \\

    \theutterance \stepcounter{utterance}  
    & \multicolumn{4}{p{0.6\linewidth}}{
        \cellcolor[rgb]{0.561,0.737,0.561}{
            \makecell[{{p{\linewidth}}}]{
                \texttt{\tiny{[P1$\rangle$GM]}}
                \texttt{```json} \\
\texttt{\{} \\
\texttt{    "current\_mode": "confirming",} \\
\texttt{    "input\_rationale": "Based on all examples, the function appears to return {-}1 if a == b, otherwise it returns the sum of the digits of the smaller absolute value. The earlier confusion came from the first example ((99, 1000), 9), where 9+9=18, but the output was 9. This suggests the sum is halved when the digit count of the larger number is even. For (99, 1000): 18//2=9 (matches). For (2, {-}2): 2+2=4 (no halving, matches). For ({-}1000, {-}1000): {-}1 (matches). This seems consistent.",} \\
\texttt{    "hypothesis": "def f(a: int, b: int) {-}> int:\textbackslash{}n    \textbackslash{}"\textbackslash{}"\textbackslash{}"Returns {-}1 if a == b, else returns the sum of digits of min(abs(a), abs(b)) halved if max(abs(a), abs(b)) has even digit count.\textbackslash{}"\textbackslash{}"\textbackslash{}"\textbackslash{}n    if a == b:\textbackslash{}n        return {-}1\textbackslash{}n    def digit\_sum(x):\textbackslash{}n        return sum(int(d) for d in str(abs(x)))\textbackslash{}n    def digit\_count(x):\textbackslash{}n        return len(str(abs(x)))\textbackslash{}n    smaller = min(abs(a), abs(b))\textbackslash{}n} \\
            }
        }
    }
    & & \\ \\

    \theutterance \stepcounter{utterance}  
    & & & \multicolumn{4}{p{0.6\linewidth}}{
        \cellcolor[rgb]{0.9,0.9,0.9}{
            \makecell[{{p{\linewidth}}}]{
                \texttt{\tiny{[P1$\langle$GM]}}
                \texttt{Game Violation: ParseError: Response is not valid JSON.} \\
\texttt{Output ONLY:} \\
\texttt{One JSON with confidence\_score, current\_mode, input\_rationale, and hypothesis keys ...} \\
\texttt{OR} \\
\texttt{SOLVE: ```python ... ```} \\
            }
        }
    }
     \\ \\

    \theutterance \stepcounter{utterance}  
    & \multicolumn{4}{p{0.6\linewidth}}{
        \cellcolor[rgb]{0.561,0.737,0.561}{
            \makecell[{{p{\linewidth}}}]{
                \texttt{\tiny{[P1$\rangle$GM]}}
                \texttt{SOLVE:} \\
\texttt{```python} \\
\texttt{def f(a: int, b: int) {-}> int:} \\
\texttt{    """Returns {-}1 if a == b, else returns the sum of digits of the smaller absolute value."""} \\
\texttt{    if a == b:} \\
\texttt{        return {-}1} \\
\texttt{    def digit\_sum(x):} \\
\texttt{        return sum(int(d) for d in str(abs(x)))} \\
\texttt{    return digit\_sum(min(a, b, key=abs))} \\
\texttt{```} \\
            }
        }
    }
    & & \\ \\

    \theutterance \stepcounter{utterance}  
    & & & \multicolumn{2}{p{0.3\linewidth}}{
        \cellcolor[rgb]{0.9,0.9,0.9}{
            \makecell[{{p{\linewidth}}}]{
                \texttt{\tiny{[GM$|$GM]}}
                \texttt{Game Verdict: LOSS} \\
            }
        }
    }
    & & \\ \\

    \theutterance \stepcounter{utterance}  
    & & & \multicolumn{2}{p{0.3\linewidth}}{
        \cellcolor[rgb]{0.9,0.9,0.9}{
            \makecell[{{p{\linewidth}}}]{
                \texttt{\tiny{[GM$|$GM]}}
                \texttt{Incorrect.} \\
\texttt{Failed Cases:} \\
\texttt{Input: {[}99, 1000{]} | Expected: 9 | Got: 18} \\
\texttt{Input: {[}{-}99, 0{]} | Expected: 9 | Got: 0} \\
\texttt{Input: {[}2, {-}2{]} | Expected: 4 | Got: 2} \\
            }
        }
    }
    & & \\ \\

    \theutterance \stepcounter{utterance}  
    & & & \multicolumn{2}{p{0.3\linewidth}}{
        \cellcolor[rgb]{0.9,0.9,0.9}{
            \makecell[{{p{\linewidth}}}]{
                \texttt{\tiny{[GM$|$GM]}}
                \texttt{def last\_digit\_sum(a: int, b: int) {-}> int:} \\
\texttt{    return abs(a) \% 10 + abs(b) \% 10} \\
            }
        }
    }
    & & \\ \\

    \theutterance \stepcounter{utterance}  
    & & & \multicolumn{2}{p{0.3\linewidth}}{
        \cellcolor[rgb]{0.9,0.9,0.9}{
            \makecell[{{p{\linewidth}}}]{
                \texttt{\tiny{[GM$|$GM]}}
                \texttt{def f(a: int, b: int) {-}> int:} \\
\texttt{    """Returns {-}1 if a == b, else returns the sum of digits of the smaller absolute value."""} \\
\texttt{    if a == b:} \\
\texttt{        return {-}1} \\
\texttt{    def digit\_sum(x):} \\
\texttt{        return sum(int(d) for d in str(abs(x)))} \\
\texttt{    return digit\_sum(min(a, b, key=abs))} \\
            }
        }
    }
    & & \\ \\

    \theutterance \stepcounter{utterance}  
    & & & \multicolumn{2}{p{0.3\linewidth}}{
        \cellcolor[rgb]{0.9,0.9,0.9}{
            \makecell[{{p{\linewidth}}}]{
                \texttt{\tiny{[GM$|$GM]}}
                \texttt{Mode: passive\_labeled\_pairs} \\
\texttt{Accuracy: 0.300} \\
\texttt{Information requests: 3/15} \\
\texttt{Revealed passive examples: 4} \\
\texttt{Efficiency: 0.867} \\
\texttt{Slack: 12} \\
\texttt{Observed{-}pair consistency: 0.000} \\
\texttt{Membership balance: 2 true / 2 false} \\
\texttt{Parse errors: 6} \\
\texttt{Runtime errors: 0} \\
            }
        }
    }
    & & \\ \\

\end{supertabular}
}

\end{document}
